% Part: incompleteness
% Chapter: representability-in-q
% Section: representing-relations

\documentclass[../../../include/open-logic-section]{subfiles}

\begin{document}

\olfileid[tr]{inc}{req}{rel}

\olsection{Bağıntıların Temsil Edilmesi}

Bir \emph{bağıntının} temsil edilebilir olmasının ne anlama geldiğini
söyleyelim.

\begin{defn}
\ollabel{defn:representing-relations} Doğal sayılar üzerindeki bir
$R(x_0,\dots,x_k)$ bağıntısı, $R(n_0,\dots,n_k)$ doğru olduğunda $\Th{Q}$
$!A_R(\num{n_0},\dots,\num{n_k})$'yi, $R(n_0,\dots,n_k)$ yanlış olduğunda ise $\Th{Q}$
$\lnot !A_R(\num{n_0},\dots,\num{n_k})$'yi ispatlayacak biçimde bir
$!A_R(x_0,\dots,x_k)$ formülü varsa {\em $\Th{Q}$ içinde temsil
edilebilir}.
\end{defn}

\begin{thm}
\ollabel{thm:representing-rels} Bir bağıntı $\Th{Q}$ içinde ancak ve ancak
hesaplanabilirse temsil edilebilir.
\end{thm}

\begin{proof}
İleri yön için $R(x_0,\dots,x_k)$'nin $!A_R(x_0,\dots,x_k)$ formülüyle
temsil edildiğini varsayın. $R$'yi hesaplayan bir algoritma şöyledir:
$n_0$, \dots,~$n_k$ girdisinde aynı anda
$!A_R(\num{n_0},\dots,\num{n_k})$'nin bir ispatını ve
$\lnot !A_R(\num{n_0},\dots,\num{n_k})$'nin bir ispatını arayın.
Varsayımımız uyarınca arama mutlaka ikisinden birini bulur; birincisini
bulursa ``evet'', ötekini bulursa ``hayır'' yanıtını verin.

Öteki yönde $R(x_0,\dots,x_k)$'nin hesaplanabilir olduğunu varsayın. Tanım
gereği bu, $\Char{R}(x_0,\dots,x_k)$ fonksiyonunun hesaplanabilir olması
demektir. \olref[int]{thm:representable-iff-comp} uyarınca $\Char{R}$ bir
formülle, diyelim $!A_{\Char{R}}(x_0,\dots,x_k,y)$ ile temsil edilir.
$!A_R(x_0,\dots,x_k)$, $!A_{\Char{R}}(x_0,\dots,x_k,\num{1})$ formülü
olsun. O zaman her $n_0$, \dots,~$n_k$ için $R(n_0,\dots,n_k)$ doğruysa
$\Char{R}(n_0,\dots,n_k)=1$ olur; bu durumda $\Th{Q}$
$!A_{\Char{R}}(\num{n_0},\dots,\num{n_k},\num{1})$'i ve dolayısıyla
$\Th{Q}$, $!A_R(\num{n_0},\dots,\num{n_k})$'yi ispatlar. Öte yandan
$R(n_0,\dots,n_k)$ yanlışsa $\Char{R}(n_0,\dots,n_k)=0$ olur. Bu,
$\Th{Q}$'nun
\[
\lforall[y][(!A_{\Char{R}}(\num{n_0}, \dots, \num{n_k}, y) \lif y =
  \num{0})].
\]
formülünü ispatladığı anlamına gelir. $\Th{Q}$,
$\eq/[\num{0}][\num{1}]$'i de ispatladığından $\Th{Q}$,
$\lnot !A_{\Char{R}}(\num{n_0},\dots,\num{n_k},\num{1})$'i ve dolayısıyla
$\lnot !A_R(\num{n_0},\dots,\num{n_k})$'yi ispatlar.
\end{proof}

\begin{prob}
$R$, $\Th{Q}$ içinde temsil edilebiliyorsa $\Char{R}$'nin de temsil
edilebildiğini gösterin.
\end{prob}

\end{document}
